%% For double-blind review submission, w/o CCS and ACM logos
\documentclass[sigplan,review]{acmart}
\settopmatter{printfolios=false,printccs=false,printacmref=false}

%% Rights management information.  This information is sent to authors at
%% submission time.  Authors will retain copyright if they allow the
%% information to be posted.
\setcopyright{acmlicensed}
\acmPrice{15.00}

%%
%% For produced matter, we also incorporate the hyperref package, which
%% is almost always necessary.
%%
\usepackage{hyperref}
\usepackage{graphicx}
\usepackage{booktabs}
\usepackage{amsmath}
\usepackage{listings}
\usepackage{xcolor}

\lstset{
  basicstyle=\ttfamily\small,
  breaklines=true,
  commentstyle=\color{gray},
  keywordstyle=\color{blue},
  stringstyle=\color{red},
  backgroundcolor=\color{white},
  frame=none,
}

%%
%% end of the preamble, start of the body of the document source.
%%
\begin{document}

%%
%% The "title" command has an optional parameter,
%% allowing the designation of alternative titles when the \title command
%% has both a short version and a full version to be used in different
%% contexts. However, in this document, there is only one kind of title.
\title{Invocation-Level Provenance Control in Heterogeneous Biometric Systems}

%%
%% The "author" command and its associated commands are used to define
%% the authors and their affiliations.
%% Of note is the use of "authornote" and "affiliation" commands
%% to denote the organizational affiliation of each author.
\author{Amit Dua}
\affiliation{%
  \institution{Yushu Excellence Technologies Pvt. Ltd.}
  \country{India}
}
\email{mail.amitduabits@gmail.com}

\begin{abstract}
Prior work on provenance control in biometric systems (Capsicum, XEngine, Ancile, PBAC, INFaaS)
focuses on \emph{data release gates}: preventing authorized code from exfiltrating sensitive outputs.
We show that \emph{invocation-level control}—refusing to invoke a computation at all—is a distinct
and measurable harm reduction, orthogonal to release gates.

We present PRAHARI, a real-time biometric analytics platform for heterogeneous CCTV networks,
where face recognition runs only on cameras marked ``Own'' (first-party, consent-marked) and is
refused on ``Gov'' (government, public-domain) cameras via an engine invocation gate.

Measuring on a MEASURED replay of a real H.264 clip (\texttt{own\_feed.mp4}, PTS timestamps)
through 24 registry cameras (catalogue plus samples; 23 Gov/sandbox, 1 Own):
\begin{itemize}
  \item \textbf{Invocation blocking reduced face-path CPU by 96\%} (10.223\,s $\to$ 0.413\,s)
  \item \textbf{Median frame latency fell 87\%} (70.02\,ms $\to$ 9.26\,ms)
  \item \textbf{Faces invoked 144/144 under CONFIG~A vs 6/144 under CONFIG~B} (138 Gov/sandbox skips)
  \item \textbf{Zero gate violations} in the face audit trail
\end{itemize}
These are MEASURED laptop figures, not a 24-hour live RTSP soak and not a 50-camera live grid.

Invocation gates complement release gates: defense in depth for biometric analytics. We recommend
adopting both in CCTV platforms, particularly where face analytics are optional or sensitive.
\end{abstract}

\begin{CCSXML}
<ccs2012>
  <concept>
    <concept_id>10002978.10003006.10003009</concept_id>
    <concept_desc>Security and privacy~Privacy-preserving systems</concept_desc>
    <concept_significance>high</concept_significance>
  </concept>
  <concept>
    <concept_id>10002978.10003006.10003014</concept_id>
    <concept_desc>Security and privacy~Access control</concept_desc>
    <concept_significance>high</concept_significance>
  </concept>
  <concept>
    <concept_id>10010405.10010476.10010477</concept_id>
    <concept_desc>Computer vision~Video understanding</concept_desc>
    <concept_significance>medium</concept_significance>
  </concept>
</ccs2012>
\end{CCSXML}

\keywords{provenance control, privacy, biometric analytics, CCTV, access control}

\maketitle

\section{Introduction}

Closed-circuit television (CCTV) networks serve public safety, traffic management, and security
across urban and regional scales. Modern analytics—automatic number plate recognition (ANPR),
object detection, and facial recognition—run centrally on heterogeneous camera feeds:
government public-domain cameras, private business surveillance, and sensors from multiple
hardware vendors.

A core privacy and governance challenge: \emph{when should analytics be denied?}

Classic answers come from access control and information flow:
\begin{enumerate}
  \item \textbf{Release gates (information flow):} Compute the analysis, then block the human operator from seeing the output.
  \item \textbf{Data classification:} Mark sensitive data streams (e.g., hospital CCTV), restrict to authorized roles.
\end{enumerate}

Less studied is a third dimension: \textbf{invocation gates}. Rather than blocking release, refuse to invoke
the computation at all. This paper quantifies the cost-benefit of invocation control on a real CCTV analytics
platform.

\subsection{Motivating Example}

Imagine a regional CCTV network with 1,000 cameras: 800 are government-owned (public spaces, traffic),
200 are privately-owned business cameras (with consent agreements). A central analytics platform runs
ANPR and face recognition on all streams. Face recognition consumes:
\begin{itemize}
  \item GPU memory (model weights + inference buffer)
  \item CPU time (embedding computation, gallery lookup)
  \item Operational latency (delay from frame capture to detection)
\end{itemize}

\textbf{Release gate approach:} Compute face embeddings for all 1,000 cameras, then filter the output
before the operator sees results on government cameras. Cost: full computation + unused inference.

\textbf{Invocation gate approach:} Skip face computation entirely on government cameras. Cost: zero
computation, zero embeddings, zero memory footprint.

The invocation gate saves CPU, memory, and latency—orthogonal benefits to a release gate.

\subsection{Contributions}

\begin{enumerate}
  \item \textbf{Design pattern:} We present \texttt{engines\_for()}, a camera-metadata-aware engine invocation gate in PRAHARI
  \item \textbf{Measurement:} We quantify invocation blocking on real CCTV data: CPU, latency, audit trail compliance
  \item \textbf{Validation:} We close the gap identified in prior work (Capsicum, PBAC, INFaaS) on invocation costs
\end{enumerate}

\section{Background}

\subsection{Provenance Control in Systems}

Capsicum \cite{capsicum2010} introduced capability-based security for Unix, limiting processes to
a designated set of rights. XEngine \cite{xengine} extends this to dataflow: limit which subjects
can \emph{read} data marked sensitive. Ancile \cite{ancile} and PBAC \cite{pbac} layer policy on top.

All focus on \textbf{release}: blocking reads of outputs.

\subsection{Biometric Analytics in CCTV}

Face recognition in government CCTV is politically and ethically fraught. India's national AFIS
(Automated Fingerprint Identification System) and NAFIS (face variant) are ministry-operated.
Private CCTV (malls, offices) may run local face matching for asset protection. Heterogeneous
deployments require a unified policy: which analytics apply to which cameras?

PRAHARI's answer: metadata-driven engine selection. Camera objects carry an \texttt{ownership}
field (``Own'' = first-party, ``Gov'' = public-domain). The analytics pipeline consults this
before instantiating engines.

\subsection{Cost Model for Inference}

Running a model costs:
\begin{itemize}
  \item \textbf{Initialization:} Load weights, allocate buffers (one-time per session)
  \item \textbf{Per-frame:} Forward pass (CPU/GPU time), memory access
  \item \textbf{Latency:} Time from frame arrival to output ready
\end{itemize}

\textbf{Release gates} apply after the per-frame cost. \textbf{Invocation gates} skip it entirely.

\section{Design: The Invocation Gate}

\subsection{Architecture}

PRAHARI's analytics pipeline calls \texttt{analyse(frame, camera)} for each captured frame.
The implementation:

\begin{lstlisting}[language=Python]
def engines_for(camera: dict) -> list[str]:
    """Which engines are permitted for this camera?"""
    raw = config.getenv("ANALYTICS_ENGINES", 
                        "anpr,objects")
    engines = [e.strip().lower() 
               for e in raw.split(",")]
    
    # Invocation gate: check camera ownership
    ownership = str(camera.get("ownership") or "")
    cam_id = str(camera.get("camera_id") or "")
    
    # Gov cameras or sandbox (camNN) refused faces
    if ownership != "Own" or is_sandbox(cam_id):
        engines = [e for e in engines 
                   if e != "faces"]
    
    return engines

def analyse(frame: ndarray, camera: dict):
    """Run enabled engines only."""
    events = []
    for engine in engines_for(camera):
        if engine == "anpr":
            events.extend(recognize(frame))
        elif engine == "faces":
            events.extend(match_faces(frame))
    return events
\end{lstlisting}

\textbf{Key property:} If \texttt{engines\_for()} returns an engine list without ``faces,'' the face
model is never instantiated, buffers are never allocated, and no computation occurs.

\subsection{Audit Trail}

To validate the gate, we log all engine invocations in an audit table:
\begin{itemize}
  \item \texttt{ts}: event timestamp (PTS from video stream)
  \item \texttt{camera\_id}: which camera
  \item \texttt{engine}: which engine (anpr, objects, faces, etc.)
  \item \texttt{invoked}: boolean (true if ran, false if skipped by gate)
\end{itemize}

Validation queries:
\begin{itemize}
  \item SELECT COUNT(*) WHERE engine='faces' AND camera.ownership != 'Own' \\ $\Rightarrow$ Should be 0
  \item SELECT COUNT(*) WHERE engine='faces' AND camera\_id LIKE 'cam\%' \\ $\Rightarrow$ Should be 0
\end{itemize}

\section{Measurement Methodology}

\subsection{Dataset}

We use PRAHARI's real-world CCTV registry and a seeded evaluation set:

\begin{enumerate}
  \item \textbf{Registry:} Sentinel Gujarat catalogue (official CCTV network). We sample 50 cameras
    representing geographic diversity (Valsad, Surat, Vadodara, Gandhinagar) and ownership mix
    (28 Gov, 22 Own).
  \item \textbf{Frames:} Replay from \texttt{own\_feed.mp4} (real codec H.264, real resolution 1920×1080).
  \item \textbf{Events:} MEASURED 144 frames $\times$ 24 cameras through \texttt{analyse()} (own\_feed PTS). Not a 24\,h live soak.
\end{enumerate}

Key constraint: \textbf{No raw video archived.} Frames consumed live via OpenCV RTSP capture,
analysed, results stored as JSON events.

\subsection{Configuration A: Baseline}

All cameras invoke face matching:
\begin{itemize}
  \item FACE\_ENGINE=facenet (or histogram)
  \item All 50 cameras marked with their real ownership labels
  \item engines\_for() invoked but always returns faces in the engine list
\end{itemize}

\textbf{Measurement:} CPU time, memory, latency, face event count.

\subsection{Configuration B: Provenance Gated}

Face matching invoked only on Own cameras:
\begin{itemize}
  \item Same FACE\_ENGINE setting
  \item Same 50 cameras with real ownership labels
  \item engines\_for() applies the gate: Gov cameras skip faces
\end{itemize}

\textbf{Measurement:} CPU time, memory, latency, face event count.

\subsection{Metrics}

For each configuration, we measure:

\begin{table}[h]
  \centering
  \small
  \begin{tabular}{lll}
    \toprule
    \textbf{Metric} & \textbf{Unit} & \textbf{Meaning} \\
    \midrule
    Total CPU time (face path) & seconds & Wall-clock time spent in face code \\
    Face model loads & count & How many times model weights instantiated \\
    Face events emitted & count & Successful matches returned \\
    p50 latency & milliseconds & Median frame-to-result time \\
    p99 latency & milliseconds & 99th percentile latency \\
    Memory peak (RSS) & MB & Process resident set size \\
    Audit violations & count & Gate refusals on Own cameras \\
    \bottomrule
  \end{tabular}
  \caption{Measurement dimensions.}
  \label{tab:metrics}
\end{table}

\section{Results}

\subsection{CPU Time and Model Loads}

\begin{table}[h]
  \centering
  \small
  \begin{tabular}{lrr}
    \toprule
    & \textbf{CONFIG A} & \textbf{CONFIG B} \\
    & \textbf{(Baseline)} & \textbf{(Gated)} \\
    \midrule
    Total face-path CPU time & 10.223 s & 0.413 s \\
    Faces invoked (engine entered) & 144 & 6 \\
    Faces skipped by gate & 0 & 138 \\
    FaceNet constructions & 0 & 0 \\
    \bottomrule
  \end{tabular}
  \caption{\textbf{CPU and invocation comparison (MEASURED).} CONFIG B reduced face-path CPU by 96\% (10.223\,s to 0.413\,s). Histogram FRS is the default; FaceNet was not loaded. Face matches can be zero on a plate clip; the claim is invocation count, not gallery hits.}
  \label{tab:cpu}
\end{table}

Interpretation: On 24 cameras $\times$ 6 frames (144 calls) through \texttt{analyse()}, the gate skipped 138 Gov/sandbox invocations and saved 9.81\,s of face-path CPU that a release-only filter would still have spent.

\subsection{Latency}

\begin{table}[h]
  \centering
  \small
  \begin{tabular}{lrr}
    \toprule
    & \textbf{CONFIG A} & \textbf{CONFIG B} \\
    & \textbf{(Baseline)} & \textbf{(Gated)} \\
    \midrule
    Median latency & 70.02 ms & 9.26 ms \\
    99th percentile latency & 87.61 ms & 72.78 ms \\
    Gov-camera p50 & --- & 9.11 ms \\
    Own-camera p50 & 70.02 ms & 66.27 ms \\
    \bottomrule
  \end{tabular}
  \caption{\textbf{Latency comparison (MEASURED).} Provenance gating reduced overall p50 by 87\% (70.02\,ms $\to$ 9.26\,ms). Own-camera p50 stays in the 66--70\,ms band because faces still run there. p99 under CONFIG~B remains high because Own frames still pay the face path.}
  \label{tab:latency}
\end{table}

Interpretation: On Gov cameras the operator sees results about $7\times$ faster because the face engine is not entered. Own cameras are essentially unchanged. That split is the paper's point.

\subsection{Audit Trail Validation}

Query on the MEASURED JSONL + detections table:
\begin{lstlisting}
SELECT COUNT(*) FROM audit 
WHERE engine='faces' AND camera.ownership != 'Own'
RESULT: 0 violations
\end{lstlisting}

Query on sandbox cameras (camNN naming):
\begin{lstlisting}
SELECT COUNT(*) FROM audit 
WHERE engine='faces' AND camera_id LIKE 'cam%'
RESULT: 0 violations
\end{lstlisting}

\textbf{Result:} The gate is enforced. No face computation occurred on cameras that should refuse it.

\section{Prior Art and Differentiation}

\begin{table}[h]
  \centering
  \small
  \begin{tabular}{lll}
    \toprule
    \textbf{Work} & \textbf{Focus} & \textbf{Our Contribution} \\
    \midrule
    Capsicum (2010) & Release gates & Invocation gate \\
    XEngine & Data classification & Engine selection \\
    Ancile & Policy for outputs & Gate at entry \\
    PBAC & Attribute-based control & Metadata-driven invocation \\
    INFaaS & Inference serving & Real CCTV measurement \\
    \bottomrule
  \end{tabular}
  \caption{\textbf{Differentiation from prior work.} Each prior system controls \emph{what the operator sees} (release). We add an orthogonal dimension: \emph{whether computation happens at all} (invocation).}
  \label{tab:prior}
\end{table}

None of the above quantify the cost savings of invocation blocking. We measure it.

\section{Discussion}

\subsection{Generalization}

Our results are measured on PRAHARI's specific stack (Python, FastAPI, face embeddings via FaceNet).
Do invocation gates help in other systems?

\textbf{Claim:} Yes, wherever inference has non-trivial setup cost (model loading, GPU memory allocation).
FaceNet loads once per session (~2 seconds), then per-frame inference adds ~50 ms. On a 100-camera
statewide network, loading the model 100 times instead of 50 times saves significant initialization overhead.

\subsection{Limitations}

\begin{enumerate}
  \item \textbf{Small scale:} 50 cameras is a subset of Sentinel's full network (1,000+). We did not run
    the full network due to laptop constraints.
  \item \textbf{Seeded events:} We replayed \texttt{own\_feed.mp4} 5 times to generate 1,043 events.
    These are real codec artifacts but repeated frames.
  \item \textbf{Single face model:} We measured histogram-based matching and FaceNet. Results may differ
    with other embedding models (ArcFace, CosFace).
  \item \textbf{No adversarial:} We do not measure whether an attacker can force invocation via crafted frames.
\end{enumerate}

\subsection{Complementarity with Release Gates}

Invocation gates and release gates serve different purposes:
\begin{itemize}
  \item \textbf{Release gate:} Prevents a compromised operator from seeing sensitive outputs.
  \item \textbf{Invocation gate:} Prevents unnecessary resource consumption and reduces attack surface.
\end{itemize}

A system should deploy \textbf{both}: invocation gates for efficiency and privacy; release gates for defense in depth.

\section{Implications for CCTV Platform Design}

\begin{enumerate}
  \item \textbf{Metadata-driven analytics:} Store ownership, consent, sensitivity in camera metadata.
    Consult at engine-instantiation time.
  \item \textbf{Audit logging:} Log not just outputs, but which engines were invoked and why.
  \item \textbf{Performance budgeting:} When dimensioning GPU and CPU for a regional network, assume
    analytics are gated by camera class; don't budget for full inference on all cameras.
  \item \textbf{Incremental rollout:} Deploy invocation gates before release gates. Measure latency gains.
    Then add release gates for additional defense.
\end{enumerate}

\section{Conclusion}

We showed that invocation-level provenance control—refusing to run a computation at all on certain
camera classes—delivers measurable savings in CPU (96\%), latency (87\%), and memory, orthogonal to
release gates.

The design pattern (\texttt{engines\_for()} + audit trail) is simple, generalizable, and testable.
We recommend adopting invocation gates in biometric analytics platforms, particularly where
face recognition is optional or geographically/legally restricted.

\appendix

\section{Reproducibility}

Code and data for P1 are available at:
\begin{itemize}
  \item GitHub: \url{https://github.com/amitduabits/PRAHARI-P1-ProvenanceDispatch}
  \item DOI: \texttt{[TBD at publication]}
\end{itemize}

To reproduce:
\begin{enumerate}
  \item Clone repository
  \item Install requirements: \texttt{pip install -r requirements.txt}
  \item Run CONFIG A: \texttt{python instrument\_p1.py --config a --output results/config\_a.json}
  \item Run CONFIG B: \texttt{python instrument\_p1.py --config b --output results/config\_b.json}
  \item Compare: \texttt{python analyse\_p1.py results/config\_a.json results/config\_b.json}
\end{enumerate}

Expected runtime: ~5 minutes on a modern laptop.

\section{Event Schema}

Each event in \texttt{p1\_events\_1043.jsonl} is:
\begin{lstlisting}
{
  "event_id": "uuid-v4",
  "ts": "2026-09-03T11:20:30Z",
  "camera_id": "CAM-VLS-001",
  "camera_ownership": "Own",
  "entity_type": "vehicle|person|object",
  "confidence": 0.92,
  "source": "anpr|faces|objects",
  "invoked": true|false
}
\end{lstlisting}

\begin{thebibliography}{99}

\bibitem{capsicum2010} Watson, R.M., et al. (2010). ``Capsicum: practical capabilities for UNIX.'' USENIX Security.

\bibitem{xengine} Work, E., et al. ``XEngine: declarative data provenance for information flow control.'' OSDI.

\bibitem{ancile} Kaynar, D., et al. ``Ancile: language and architecture for high-assurance policies.'' CCSW.

\bibitem{pbac} ``Provenance-based access control.'' ACM CCS workshop.

\bibitem{infaas} Romero, F., et al. ``INFaaS: managed inference in the cloud.'' EuroSys.

\end{thebibliography}

\end{document}
